%% Bioinformatics.cls v2.0
%% A LaTeX class file for Oxford University Press Bioinformatics Journal
%% App Note Template

\documentclass{bioinfo}
\copyrightyear{2026} \pubyear{2026}

\access{Advance Access Publication Date: Day Month Year}
\appnotes{Application Note}

\begin{document}
\firstpage{1}

\title[Skeptic Engine]{Skeptic Engine: Calibrated, Adversarial Anomaly Detection for Scientific Data Integrity}

\author[Boiko]{Sergey~V.~Boiko\,$^{1,*}$}
\address{$^{1}$Independent Researcher, Scientific Data Integrity Group}

\corresp{$^\ast$To whom correspondence should be addressed. E-mail: sergeikuch80@gmail.com}

\history{Received on XXXXX; revised on XXXXX; accepted on XXXXX}

\editor{Associate Editor: XXXXXXX}

\abstract{\textbf{Motivation:} As the "Replication Crisis" highlights the fragility of scientific claims, the need for automated tools to screen for statistical artifacts and data integrity issues grows. While specialized tools exist for p-hacking (p-curve) or image forensics, no unified framework exists to check the structural coherence of high-dimensional scientific data across modalities.\\
\textbf{Results:} We present \textbf{Skeptic Engine}, an adversarial testing framework that transfers anomaly detection methods from finance and clinical trials to science. Our system integrates multi-modal detectors (Benford forensics, autoencoders, behavioral patterns) into a \textbf{Unified Anomaly Score (UAS)}. Crucially, we introduce \textbf{Isotonic Recalibration} to transform raw scores into calibrated probabilities with confidence intervals (reducing MACE from 0.202 to 0.032). Evaluated across 37 experiments, our approach detects statistical artifacts with AUC 0.729--1.000. The system also employs an \textbf{Adversarial Debate Protocol} to generate interpretable verdicts with evidence trails.\\
\textbf{Availability:} The open-source implementation and documentation are available at \url{https://github.com/sergeeey/Skeptic-Engine-2026-}.\\
\textbf{Supplementary information:} Supplementary data are available at \textit{Bioinformatics} online.}

\maketitle

\section{Introduction}
Statistical integrity tools often operate in isolation. P-curve analysis \citep{simonsohn2014p} examines p-value distributions but ignores data structure. Statcheck \citep{nuijten2016prevalence} recomputes test statistics but does not analyze raw data matrices. Image forens tools detect pixel manipulation but are silent on statistical coherence. Each answers a narrow question: "Are the p-values suspicious?" None asks the broader question: "Does this dataset exhibit structural coherence consistent with genuine measurement?"

Furthermore, existing detectors rarely provide calibrated uncertainty. A score of "0.8" is meaningless without knowing the probability of a true negative scoring 0.8. Without calibration, reviewers cannot distinguish between a borderline flag and a confident detection.

We propose a shift from Single-Point Detection to \textbf{Calibrated, Adversarial Verification}. \textbf{Skeptic Engine} integrates seven specialized detectors into a Unified Anomaly Score (UAS), calibrates raw scores into interpretable probabilities, and synthesizes findings through an adversarial debate protocol.

\section{Features}
\subsection{Multi-Modal Detection}
The system operates in four stages:
\begin{enumerate}
    \item \textbf{Extraction:} Features from raw data (Benford digits, p-value sequences, cross-modal correlations).
    \item \textbf{Detection:} Specialized detectors (H24--H33) compute raw anomaly scores.
    \item \textbf{Calibration (H34):} Isotonic Regression maps scores to calibrated probabilities.
    \item \textbf{Verdict (H36):} An adversarial Debate Protocol generates a structured decision with explanation.
\end{enumerate}

\subsection{Unified Anomaly Score (UAS)}
We compute a weighted ensemble: $UAS = \sum w_i \cdot S_i$. By combining signals, UAS captures complementary error profiles. Benford analysis detects digit-level anomalies that autoencoders miss; autoencoders catch structural incoherence that Benford analysis misses.

\subsection{Isotonic Recalibration}
To address the "Black Box" problem, we train a non-decreasing isotonic regression model $f_{iso}$ on historical ground-truth data: $P(\text{Anomaly} | x) \approx f_{iso}(\text{Raw Score}(x))$. This yields calibrated probabilities with Confidence Intervals based on the local density of calibration data (16,224 samples used).

\section{Results}
We evaluate across four data modalities: scRNA-seq (PBMC3k, Kang2018), Proteomics/CNA (CPTAC), P-value sequences (RPP, Statcheck), and Cross-omics (matched mRNA-protein).

\begin{table}[t]
\caption{Performance of individual detectors vs. Unified Anomaly Score}
\begin{tabular}{llc}
\hline
\textbf{Experiment} & \textbf{Domain} & \textbf{AUC} \\
\hline
H24: Benford Forensics & scRNA-seq & 0.978 \\
H25: Banking Autoencoder & Proteomics/CNA & 1.000 \\
H23: Behavioral P-Hacking & Meta-analysis & 0.729 \\
\hline
\textbf{H31: Unified Anomaly Score} & \textbf{Multi-Modal} & \textbf{1.000} \\
\hline
\end{tabular}
\end{table}

\subsection{Calibration Performance}
Recalibration significantly improves interpretability. We measured the Mean Absolute Calibration Error (MACE) before and after isotonic regression:
\begin{itemize}
    \item \textbf{Baseline MACE:} 0.202
    \item \textbf{Calibrated MACE:} 0.032 (84\% improvement)
\end{itemize}

\subsection{Adversarial Verdicts}
The Debate Protocol provides explainability. In a test case with mixed signals (High Benford deviation but high Cross-Modal consistency), the system rendered:
\begin{quote}
\textit{"Prosecutor: 3 anomalies found (Benford, temporal drift). Defense: Large sample size (n=5000), no batch effects. Verdict: SUSPICIOUS (Confidence 0.72)."}
\end{quote}

\section{Conclusion}
Skeptic Engine demonstrates that we can move beyond simple statistical checks to a holistic, calibrated, and explainable system for data integrity. By combining multi-modal detection with adversarial debate, we provide a tool that doesn't just find errors, but helps reviewers understand them.

\section*{Funding}
This research received no specific grant from any funding agency.

\section*{Conflict of Interest}
none declared.

\bibliographystyle{plainnat}
\bibliography{references}

\end{document}
